\documentclass[11pt]{article}

\usepackage[margin=1in]{geometry}
\usepackage{amsmath,amssymb,amsthm}
\usepackage{graphicx}
\usepackage[hidelinks]{hyperref}

\newtheorem{theorem}{Theorem}
\newtheorem{lemma}{Lemma}

\newcommand{\ellTwo}{\ell_2}
\newcommand{\N}{\mathbb N}
\newcommand{\R}{\mathbb R}

\title{A Uniformly Convex Set with an Extreme Point That Is Not Strongly Exposed}
\author{FA/Banach run fa\_banach\_001}
\date{June 14, 2026}

\begin{document}
\maketitle

\section*{Status}

This packet gives a counterexample to the question in Remark 4.15 of
Garcia-Lirola--Grelier--Rueda Zoca \cite{GGRZ2021}: in general, an extreme
point of a uniformly convex set need not be strongly exposed by an ambient
continuous functional.

\section*{Source Question}

The source paper asks whether every extreme point is a strongly exposed point
in a uniformly convex set \(K\). The authors then prove that every such point
is intrinsically strongly exposed by a functional continuous for the
Minkowski functional \(|\cdot|_K\), while noting that this functional need not
be continuous for the ambient norm.

\begin{figure}[h]
\centering
\includegraphics[width=\textwidth]{figures/open_problem_crop.png}
\caption{Source evidence from arXiv:2109.01393, PDF page 26, Remark 4.15.}
\end{figure}

\section*{Definitions Used}

We use the source paper's definition: a symmetric bounded closed convex set
\(K\) in a Banach space \(X\) is uniformly convex if for every
\(\varepsilon>0\) there is \(\delta>0\) such that, whenever \(x,y\in K\) and
\(\|x-y\|>\varepsilon\), one has
\[
  \frac{x+y}{2}\in (1-\delta)K.
\]
A strongly exposed point must in particular be exposed by some functional in
the ambient dual \(X^*\).

\section*{Proof Intuition}

The source remark identifies the possible obstruction: the Minkowski
functional of \(K\) may define a stronger Banach norm on \(\operatorname{span}
K\) than the ambient norm inherited from \(X\). The counterexample makes this
gap explicit. Take a Hilbert ball for a stronger weighted \(\ell_2\)-norm and
view it inside ordinary \(\ell_2\). Its boundary points are extreme and the set
is uniformly convex, but some Hilbert norming functionals are not continuous
for the ambient \(\ell_2\)-norm. Such a point is intrinsically strongly
exposed, exactly as in the source remark, but not exposed in the original
Banach space.

\begin{theorem}
There is a uniformly convex set \(K\) in a Banach space \(X\) and a point
\(x_0\in \operatorname{ext}(K)\) such that \(x_0\) is not strongly exposed.
\end{theorem}

\begin{proof}
Work over the real scalar field and let \(X=\ellTwo\) with its usual norm.
Define
\[
  K=\left\{x=(x_n)_{n\ge 1}\in \ellTwo:
    \sum_{n=1}^{\infty} n^2 x_n^2\le 1\right\}.
\]
This set is symmetric, bounded, and convex. It is closed in \(\ellTwo\):
if \(x^{(m)}\to x\) in \(\ellTwo\) and
\(\sum n^2 (x_n^{(m)})^2\le 1\) for all \(m\), then for every \(N\),
\[
  \sum_{n=1}^N n^2x_n^2
   =\lim_{m\to\infty}\sum_{n=1}^N n^2 (x_n^{(m)})^2
   \le 1,
\]
and letting \(N\to\infty\) gives \(x\in K\).

Let
\[
  |x|_K=\left(\sum_{n=1}^{\infty}n^2x_n^2\right)^{1/2}
\]
on \(\operatorname{span}K\). This is a Hilbert norm and \(K\) is its unit ball.
Since \(\|x\|_2\le |x|_K\), if \(u,v\in K\) and \(\|u-v\|_2>\varepsilon\), then
\(|u-v|_K>\varepsilon\). The unit ball of a Hilbert space is uniformly convex,
so there is \(\delta>0\) such that
\[
  \left|\frac{u+v}{2}\right|_K\le 1-\delta.
\]
Equivalently, \((u+v)/2\in (1-\delta)K\). Hence \(K\) is uniformly convex as a
subset of the ambient Banach space \(X=\ellTwo\).

Choose
\[
  c=\left(\sum_{n=1}^{\infty}\frac1{n^2}\right)^{-1/2},
  \qquad
  x_0=\left(\frac{c}{n^2}\right)_{n\ge 1}.
\]
Then
\[
  |x_0|_K^2=\sum_{n=1}^{\infty}n^2\frac{c^2}{n^4}
  =c^2\sum_{n=1}^{\infty}\frac1{n^2}=1.
\]
Because \(K\) is the unit ball of a strictly convex Hilbert norm, every
\(|\cdot|_K\)-unit vector is an extreme point of \(K\). Thus
\(x_0\in\operatorname{ext}(K)\).

It remains to prove that \(x_0\) is not exposed by any functional in
\(X^*=\ellTwo\). Let \(b=(b_n)\in\ellTwo\), and let
\[
  \varphi_b(x)=\sum_{n=1}^{\infty}b_nx_n
\]
be the corresponding ambient continuous functional. On the Hilbert space
\((\operatorname{span}K,|\cdot|_K)\), this functional is represented by the
vector
\[
  z=\left(\frac{b_n}{n^2}\right)_{n\ge 1},
  \qquad
  \varphi_b(y)=\sum_{n=1}^{\infty}n^2 z_n y_n
  =\langle y,z\rangle_K .
\]
If \(\varphi_b\) exposed \(x_0\) over \(K\), then in particular it would attain
its maximum over the Hilbert unit ball \(K\) at \(x_0\). Equality in the
Cauchy--Schwarz inequality for the Hilbert norm forces
\[
  z=\lambda x_0
\]
for some \(\lambda>0\). Therefore
\[
  b_n=n^2z_n=\lambda n^2x_{0,n}=\lambda c
\]
for every \(n\). This coefficient sequence is constant and nonzero, hence it
does not belong to \(\ellTwo\), a contradiction.

So \(x_0\) is not exposed by any ambient continuous functional. Since every
strongly exposed point is exposed, \(x_0\) is not strongly exposed. This
finishes the counterexample.
\end{proof}

\section*{Verification Notes}

The proof uses no computational verifier. The checks are structural:
\begin{itemize}
  \item \(K\) is closed by finite-coordinate truncation and monotone
  convergence.
  \item Uniform convexity follows from the Hilbert unit ball for the stronger
  norm \(|\cdot|_K\) and the inequality \(\|x\|_2\le |x|_K\).
  \item The chosen point \(x_0=(c/n^2)\) lies on the \(|\cdot|_K\)-unit sphere,
  hence is extreme.
  \item Ambient exposing functionals correspond to coefficient sequences in
  \(\ell_2\). The only possible Hilbert norming coefficient sequence for
  \(x_0\) is constant, hence not in \(\ell_2\).
\end{itemize}

\section*{Novelty and Search Bounds}

Before packaging, the run's cheap indexes were searched for the arXiv id and
the core terms ``uniformly convex set'', ``strongly exposed point'', and
``extreme point''. The local parsed arXiv corpus was searched for close
variants including ``uniformly convex set strongly exposed'', ``extreme point
not strongly exposed'', ``compact ellipsoid'', and ``Minkowski functional
strongly exposed''. A web search on June 14, 2026 for the exact question and
close variants surfaced the source paper itself as the only relevant hit. No
prior run packet or direct literature answer for this counterexample was
found.

\section*{Human Review Recommendation}

Review as a short counterexample packet. The key point for verification is the
no-exposure argument: if an ambient \(\ell_2\) functional exposed \(x_0\), its
weighted-Hilbert representative would have to be a positive scalar multiple of
\(x_0\), forcing the ambient coefficient sequence to be constant and nonzero.

\begin{thebibliography}{9}

\bibitem{GGRZ2021}
L. C. Garcia-Lirola, G. Grelier, and A. Rueda Zoca,
\emph{Extremal structure in ultrapowers of Banach spaces},
arXiv:2109.01393.

\end{thebibliography}

\end{document}
